% Generated mechanically from the reviewed Markdown source.
% Controller edits should be made in both source and LaTeX when material.

\section{Case Study: Tracing an Animal-Navigation Deck}

The formal ThreadBench reference case asks for a science-teaching deck on animal navigation, coupling geomagnetic, solar, and stellar compasses with mechanism explanations, scientific diagrams, real images, and an exact six-column-by-three-row comparison table. This creates one long-range thread: evidence established during research must survive page planning and mechanism exposition, then reappear under a consistent comparison schema. Research records the evidence and terminology; the Orchestrator maps each requirement to deck and page contracts plus image tasks; Group Agents jointly author related mechanism and synthesis pages; and whole-deck Review checks the closing comparison against the concepts introduced earlier. The persistent artifacts expose three inspectable handoffs: evidence to plan, plan and assets to pages, and page groups to the assembled deck.

The case-specific rubric then localizes failures that look similar in the final artifact: missing mechanism evidence belongs to Research, grounded facts dropped from page contracts to Planning, and unused assets or uninspected renders to production; incorrect comparison dimensions, inconsistent terminology, and unreadable encodings remain final-deck failures even when intermediate artifacts exist. This separation gives practitioners a concrete repair target and prevents process completion from substituting for artifact quality. We use the case as a qualitative trace, while the controlled 9B ablations provide the causal test of responsibility topology. The evidence boundary remains narrow: 9B and 27B differ in data and training, LLM judgments remain bounded by their measured human agreement, the benchmark currently has one formal case, and HTML-to-PDF/PPTX conversion does not test speech, rehearsal, animation, or native-office editing. The case therefore supports failure localization, not broad domain-generalization claims.

\section{Conclusion}

Long-horizon presentation authoring fails when decisions that persist across pages outlive the context or owner responsible for them. MuralPresenter addresses this mismatch by externalizing deck state, assigning dependent pages to shared Group Agents, and combining group-level pixel inspection with independent whole-deck review. We train MuralPresenter-9B on 1,000 quality-controlled trajectories and MuralPresenter-27B on a larger heterogeneous cohort, with both systems preserving planning, diagnosis, and repair supervision. Across PresentBench, SlidesGen-Bench, and DECKBench, the models achieve \resulttbd{[TBD:CONC\_GENERAL]}; on ThreadBench, they achieve \resulttbd{[TBD:CONC\_THREAD]}. Controlled 9B ablations show \resulttbd{[TBD:CONC\_ABLATION]}, clarifying the roles of responsibility topology, delegation, visual critique, and trajectory supervision. The animal-navigation case further traces one requirement from evidence through planning and assets into rendered pages. Limitations remain: the two model sizes are not a controlled scaling comparison, ThreadBench currently contains one formal case, and HTML-to-PDF/PPTX conversion does not evaluate rehearsal, animation, or native-office editing. Within these boundaries, the results establish MuralPresenter as an auditable framework for studying and improving long-horizon presentation consistency.
